\section{Navier--Stokes Coordinate and Energy Details}
\label{app:ns-coordinate-details}
This appendix collects mathematical orientation material for the 3D
Navier--Stokes model. It is background for the continuum hierarchy; it is not
evidence for the 1D study specification or any downstream validation claim.
All formulas are read under the standing regularity and differential-operator
conventions in Appendix~\ref{app:mathematical-foundations}. In particular,
$\Divbf$ is row-divergence for tensor fields, $\del\vu$ is the
componentwise vector Laplacian, tensor norms are Frobenius norms when written,
and $\eta$ is dynamic viscosity while $\nu=\eta/\rho$ is kinematic
viscosity as in Remark~\ref{rmk:viscosity-notation}.

\subsection{Viscous-Term Simplification and Energy Context}
\label{subsec:viscous-term-simplification-energy-context}

For a Newtonian fluid with constant dynamic viscosity $\eta$ and sufficiently
smooth velocity field,
\begin{flalign*}
    \quad \Divbf(2\eta\vct{D}(\vu)) = \eta\Divbf(\nabla\vu+\nabla\vu^\top)
    = \eta\Delta\vu+\eta\nabla(\Div\vu). &&
\end{flalign*}
For incompressible flow, $\Div\vu=0$, so
\begin{flalign*}
    \quad \Divbf(2\eta\vct{D}(\vu))=\eta\Delta\vu, \qquad
    \rho^{-1}\Divbf(2\eta\vct{D}(\vu))=\nu\Delta\vu,
    \qquad \nu=\eta/\rho. &&
\end{flalign*}
This term is the momentum-diffusion contribution; it does not disappear in the
incompressible Newtonian limit.

Classical local well-posedness results for Navier--Stokes depend on the domain,
boundary conditions, forcing, initial-data space, and solution class. Under
sufficiently smooth compatible data one obtains a unique strong solution on a
maximal time interval. If the maximal interval is finite, continuation fails
through blow-up of the norm controlled by the chosen well-posedness theory.

\subsection{Cylindrical Coordinate Form}
\label{app:cylindrical-coords}

For $(r,\theta,z)$ with orthonormal basis
$(\vct{e}_r,\vct{e}_\theta,\vct{e}_z)$,
\begin{flalign*}
    \quad x=r\cos\theta,\qquad y=r\sin\theta,\qquad
    \vu=u_r\vct{e}_r+u_\theta\vct{e}_\theta+u_z\vct{e}_z. &&
\end{flalign*}
The component formulas are local cylindrical-coordinate formulas on patches
where $r>0$. Axis regularity, wall traces, and boundary conditions require
separate hypotheses and are not supplied by the coordinate notation alone.
The differential operators used in cylindrical divided-by-density
Navier--Stokes forms are
\begin{flalign*}
    \quad (\vu\cdot\nabla) &= u_r \pd_r+\frac{u_\theta}{r} \pd_\theta+u_z \pd_z, \\ \nabla p
    &= \pd_r p\,\vct{e}_r +\frac{1}{r} \pd_\theta p\,\vct{e}_\theta +\pd_z p\,\vct{e}_z, \\ \Div\vu
    &= \frac{1}{r} \pd_r (ru_r) +\frac{1}{r} \pd_\theta u_\theta +\pd_z u_z. &&
\end{flalign*}
The rate-of-deformation components in the cylindrical orthonormal basis are
\begin{flalign*}
    \quad D_{rr} &= \pd_r u_r, &
    D_{\theta\theta} &= \frac{1}{r} \pd_\theta u_\theta+\frac{u_r}{r}, &
    D_{zz} &= \pd_z u_z, \\
    D_{r\theta} &= \frac{1}{2}\left(\pd_r u_\theta-\frac{u_\theta}{r} +\frac{1}{r} \pd_\theta u_r\right), &
    D_{rz} &= \frac{1}{2}(\pd_r u_z+\pd_z u_r), &
    D_{\theta z} &=
    \frac{1}{2}\left(\frac{1}{r} \pd_\theta u_z+\pd_z u_\theta\right). &&
\end{flalign*}
Use the scalar shear-rate convention
$\dot\gamma=\sqrt{2\vct{D}(\vu):\vct{D}(\vu)}$. Then
\begin{flalign*}
    \quad \tau_{ij}=2\eta_{\mathrm{eff}}(\dot\gamma)D_{ij}, \qquad i,j\in\{r,\theta,z\}. &&
\end{flalign*}
When a source writes the law as a function of $D:D$ or $|D|_{\mathrm F}^2$,
this report translates it through this convention, consistent with
Remark~\ref{rmk:viscosity-notation}.
Writing $\vf=f_r\vct{e}_r+f_\theta\vct{e}_\theta+f_z\vct{e}_z$, the momentum
components are
\begin{flalign*}
    \text{(radial)}\quad &\pd_t u_r
    +u_r \pd_r u_r +\frac{u_\theta}{r} \pd_\theta u_r +u_z \pd_z u_r -\frac{u_\theta^2}{r} \\
    &\quad = -\frac{1}{\rho} \pd_r p +\frac{1}{\rho}
    \left[
        \frac{1}{r} \pd_r (r\tau_{rr})
        +\frac{1}{r} \pd_\theta \tau_{r\theta}
        +\pd_z \tau_{rz}
        -\frac{1}{r}\tau_{\theta\theta}
    \right]
    +f_r, \\
    \text{(azimuthal)}\quad &\pd_t u_\theta+u_r \pd_r u_\theta
    +\frac{u_\theta}{r} \pd_\theta u_\theta+u_z \pd_z u_\theta+\frac{u_ru_\theta}{r} \\
    &\quad = -\frac{1}{\rho r} \pd_\theta p+\frac{1}{\rho}
    \left[
        \frac{1}{r} \pd_r (r\tau_{r\theta})
        +\frac{1}{r} \pd_\theta \tau_{\theta\theta}
        +\pd_z \tau_{\theta z}
        +\frac{1}{r}\tau_{r\theta}
    \right]+f_\theta, \\
    \text{(axial)}\quad &\pd_t u_z +u_r \pd_r u_z+\frac{u_\theta}{r} \pd_\theta u_z+u_z \pd_z u_z \\
    &\quad = -\frac{1}{\rho} \pd_z p+\frac{1}{\rho}
    \left[
        \frac{1}{r} \pd_r (r\tau_{rz})
        +\frac{1}{r} \pd_\theta \tau_{\theta z}
        +\pd_z \tau_{zz}
    \right]+f_z. &&
\end{flalign*}
The incompressibility condition is
\begin{flalign*}
    \quad \frac{1}{r} \pd_r (ru_r) +\frac{1}{r} \pd_\theta u_\theta+\pd_z u_z=0. &&
\end{flalign*}
